\chapter{समकोण त्रिभुज और कोज्या}\label{ch:g8:cosine}

दो सुंदर बातें \omterm{def:g6:shapes:triangles}{समकोण त्रिभुज} को वृत्त से बाँध देती हैं: \omterm{def:g6:shapes:triangles}{समकोण
त्रिभुज} ठीक आधे वृत्त में बैठ जाता है, और उसका कर्ण एक \omterm{def:g4:geometry:circle}{व्यास} होता
है। और एक अकेली संख्या, किसी कोण की \emph{कोज्या}, उस कोण वाले हर
\omterm{def:g6:shapes:triangles}{समकोण त्रिभुज} की शक्ल अपने भीतर समेट लेती है --- यह पहला
त्रिकोणमितीय अनुपात है, अपने भाई-बहनों ज्या और स्पर्शज्या से पहले
(\cref{ch:g9:trig})।

\section{समकोण त्रिभुज और उसका वृत्त}

\begin{theorem}[वृत्त प्रमेय]\label{thm:g8:cosine:circle}
\begin{enumerate}
  \item अगर किसी त्रिभुज $ABC$ का $A$ पर \omterm{def:g6:shapes:triangles}{समकोण} हो, तो $A$ उस वृत्त
        पर पड़ता है जिसका \omterm{def:g4:geometry:circle}{व्यास} कर्ण $[BC]$ है।
  \item और उलटे तरफ़ से: अगर $A$ किसी ऐसे वृत्त पर पड़े जिसका \omterm{def:g4:geometry:circle}{व्यास}
        $[BC]$ हो (और $A \neq B, C$), तो त्रिभुज $ABC$ का $A$ पर
        \omterm{def:g6:shapes:triangles}{समकोण} होता है।
\end{enumerate}
यही बात दूसरे शब्दों में: किसी \omterm{def:g6:shapes:triangles}{समकोण त्रिभुज} में कर्ण का मध्यबिंदु
तीनों \omterm{def:g5:solids:def}{शीर्षों} से बराबर दूरी पर होता है --- यानी \omterm{def:g6:shapes:triangles}{समकोण} से खींची
माध्यिका कर्ण की आधी नापती है।
\end{theorem}

\begin{proof}[1 का प्रमाण]
मान लो $O$, $[BC]$ का मध्यबिंदु है और $D$, $O$ के बारे में $A$ का
\omterm{def:g7:central:def}{सममित} बिंदु। $ABDC$ के विकर्ण अपने साझे मध्यबिंदु $O$ पर कटते हैं,
इसलिए $ABDC$ एक \omterm{def:g7:central:parallelogram}{समांतर चतुर्भुज} है
(\cref{def:g7:central:parallelogram}) --- और उसका $A$ पर \omterm{def:g6:shapes:triangles}{समकोण} है:
यानी वह आयत है। आयत के विकर्ण बराबर होते हैं, इसलिए
$OA = \frac{AD}{2} = \frac{BC}{2}$: यानी बिंदु $A$, $O$ से
$\frac{BC}{2}$ दूरी पर है, यानी \omterm{def:g4:geometry:circle}{व्यास} $[BC]$ वाले वृत्त पर। (बात 2
इसी तर्क को उलटा चलाकर सिद्ध होती है।)
\end{proof}

\begin{omfigure}
\begin{tikzpicture}[scale=1.15]
  \coordinate (B) at (-1.8,0);
  \coordinate (C) at (1.8,0);
  \coordinate (O) at (0,0);
  \coordinate (A) at (130:1.8);
  \coordinate (A2) at (55:1.8);
  \draw[thick] (O) circle (1.8);
  \draw[thick] (B) -- (C);
  \draw[thick, omDef] (B) -- (A) -- (C);
  \draw[thick, omExo] (B) -- (A2) -- (C);
  \draw[omThm, thick] (O) -- (A);
  \fill (B) circle (1.5pt) node[below left, font=\small] {$B$};
  \fill (C) circle (1.5pt) node[below right, font=\small] {$C$};
  \fill (O) circle (1.5pt) node[below, font=\small] {$O$};
  \fill (A) circle (1.5pt) node[above left, font=\small] {$A$};
  \fill (A2) circle (1.5pt) node[above right, font=\small] {$A'$};
  \draw[thin] (A) ++(0.14,-0.20) -- ++(0.14,0.10) -- ++(-0.14,0.20);
\end{tikzpicture}

{\small $A$ वृत्त पर चाहे कहीं भी हो, कोण $\widehat{BAC}$ \omterm{def:g6:shapes:triangles}{समकोण} ही
रहता है --- \cref{exo:g6:shapes:10} वाली अटकल अब प्रमेय बन गई। लाल
माध्यिका $[OA]$ एक \omterm{def:g3:shapes:shapes}{त्रिज्या} है: यानी कर्ण की आधी।}
\end{omfigure}

\begin{example}\label{ex:g8:cosine:circle}
किसी त्रिभुज का कर्ण $10$ cm है। और कुछ जाने बिना ही: \omterm{def:g6:shapes:triangles}{समकोण} से खींची
माध्यिका $5$ cm नापती है, और त्रिभुज के \omterm{pb:g6:symmetry:1}{परिवृत्त} की \omterm{def:g3:shapes:shapes}{त्रिज्या} $5$ cm
है, जिसका केंद्र कर्ण का मध्यबिंदु है।
\end{example}

\section{किसी न्यून कोण की कोज्या}

\begin{definition}[कोज्या]\label{def:g8:cosine:def}
किसी \omterm{def:g6:shapes:triangles}{समकोण त्रिभुज} में, किसी न्यून कोण $\theta$ के लिए:
\[
\cos\theta = \frac{\text{संलग्न भुजा}}{\text{कर्ण}}
\]
--- यानी $\theta$ को छूती \omterm{def:g6:shapes:triangles}{समकोण} भुजा, कर्ण से भाग दी हुई। यह अनुपात
सिर्फ़ कोण पर टिका है, त्रिभुज की नाप पर नहीं: एक ही न्यून कोण वाले
सारे \omterm{def:g6:shapes:triangles}{समकोण त्रिभुज} एक-दूसरे के बढ़े हुए रूप हैं, और बढ़ाने से
लंबाइयों के अनुपात नहीं बदलते (\cref{ch:g8:midpoints} ने यह कहानी
शुरू की थी; \cref{ch:g9:thales} उसे पूरा करता है)।
\end{definition}

\begin{omfigure}
\begin{tikzpicture}[scale=1.05]
  \coordinate (B) at (0,0);
  \coordinate (A) at (4.6,0);
  \coordinate (C) at (4.6,2.3);
  \coordinate (A1) at (2.76,0);
  \coordinate (C1) at (2.76,1.38);
  \draw[thick] (B) -- (A) -- (C) -- cycle;
  \draw[omThm, thick] (A1) -- (C1);
  \draw[thin] (A) ++(-0.25,0) -- ++(0,0.25) -- ++(0.25,0);
  \draw[thin] (A1) ++(-0.22,0) -- ++(0,0.22) -- ++(0.22,0);
  \draw[omDef, ->] (0.85,0) arc (0:26.6:0.85);
  \node[omDef, font=\small] at (1.25,0.24) {$\theta$};
  \node[below left, font=\small] at (B) {$B$};
\end{tikzpicture}

{\small एक ही कोण $\theta$ वाले दो \omterm{def:g6:shapes:triangles}{समकोण त्रिभुज}: छोटा वाला बड़े का
घटा हुआ रूप है, इसलिए
$\frac{\text{संलग्न}}{\text{कर्ण}}$ दोनों में एक ही रहता है ---
वही साझा मान $\cos\theta$ है।}
\end{omfigure}

\begin{proposition}[पहले मान और सीमाएँ]\label{prop:g8:cosine:bounds}
हर न्यून कोण $\theta$ के लिए $0 < \cos\theta < 1$ होता है, और कोण
खुलने के साथ कोज्या \emph{घटती} जाती है: बड़ा कोण, छोटी कोज्या।
याद रखने लायक़ पड़ाव: $\cos 0^\circ = 1$,
$\cos 60^\circ = \frac12$, $\cos 90^\circ = 0$ (ये सिरे वाले मान
चपटे पड़ चुके त्रिभुजों के हैं)।
\end{proposition}
\admitted

\begin{method}[कोज्या का इस्तेमाल]\label{met:g8:cosine:use}
किसी \omterm{def:g6:shapes:triangles}{समकोण त्रिभुज} में, जब पता और पूछी गई राशियाँ कोई न्यून कोण,
उसकी संलग्न भुजा और कर्ण हों:
\begin{enumerate}
  \item परिभाषा लिखो:
        $\cos\theta = \frac{\text{संलग्न}}{\text{कर्ण}}$, और उसमें
        पता मान रख दो;
  \item इस छोटे \omterm{def:g8:equations:def}{समीकरण} को अज्ञात के लिए हल करो (गुणा या भाग करके);
  \item अगर अज्ञात कोई \emph{कोण} हो, तो निकाले हुए अनुपात पर
        कैलकुलेटर का $\cos^{-1}$ लगाओ;
  \item मोटी जाँच: कोई लंबाई कर्ण से छोटी ही निकलनी चाहिए; और कोई
        कोण पक्के तौर पर $0^\circ$ तथा $90^\circ$ के बीच।
\end{enumerate}
\end{method}

\begin{example}[कोई भुजा निकालना]\label{ex:g8:cosine:side}
$6$ m का एक ढलान ज़मीन के साथ $20^\circ$ का कोण बनाता है। उससे तय
हुई आड़ी दूरी ($20^\circ$ की संलग्न, कर्ण $6$ m):
\[
\cos 20^\circ = \frac{d}{6}
\quad\Longrightarrow\quad
d = 6 \cos 20^\circ \approx 6 \times 0.940 \approx 5.6 \text{ m}.
\]
\end{example}

\begin{example}[कोई कोण निकालना]\label{ex:g8:cosine:angle}
किसी \omterm{def:g6:shapes:triangles}{समकोण त्रिभुज} में कोण $\theta$ की संलग्न भुजा $3.5$ नापती है
और कर्ण $5$:
\[
\cos\theta = \frac{3.5}{5} = 0.7,
\qquad
\theta = \cos^{-1}(0.7) \approx 45.6^\circ .
\]
\end{example}

\section{अभ्यास}

\begin{exercise}[$\star$]\label{exo:g8:cosine:1}
किसी \omterm{def:g6:shapes:triangles}{समकोण त्रिभुज} का कर्ण $13$ cm है। \omterm{def:g6:shapes:triangles}{समकोण} वाले \omterm{def:g5:solids:def}{शीर्ष} से खींची
माध्यिका कितनी लंबी है? और तीनों \omterm{def:g5:solids:def}{शीर्षों} से गुज़रते वृत्त की
\omterm{def:g3:shapes:shapes}{त्रिज्या} कितनी है?
\end{exercise}

\begin{exercise}[$\star$]\label{exo:g8:cosine:2}
$8$ cm \omterm{def:g4:geometry:circle}{व्यास} का एक वृत्त बनाओ जिसमें $[BC]$ एक \omterm{def:g4:geometry:circle}{व्यास} हो, वृत्त पर
कोई भी बिंदु $A$ चुनो और $ABC$ बनाओ। कौन-सा कोण \omterm{def:g6:shapes:triangles}{समकोण} है? प्रमेय का
नाम लो। $AB$ और $AC$ नापकर पाइथागोरस जाँचो।
\end{exercise}

\begin{exercise}[$\star$]\label{exo:g8:cosine:3}
$D$ पर \omterm{def:g6:shapes:triangles}{समकोण} वाले त्रिभुज $DEF$ में, कोण $\widehat E$ की संलग्न
भुजा, उसके सामने वाली भुजा, और कर्ण के नाम बताओ। फिर
$\cos \widehat E$ एक अनुपात के रूप में लिखो।
\end{exercise}

\begin{exercise}[$\star$]\label{exo:g8:cosine:4}
छूटी हुई राशि निकालो ($\cos 35^\circ \approx 0.819$,
$\cos 50^\circ \approx 0.643$):
\begin{enumerate}
  \item कर्ण $10$, कोण $35^\circ$: संलग्न भुजा?
  \item संलग्न भुजा $6$, कोण $50^\circ$: कर्ण?
\end{enumerate}
\end{exercise}

\begin{exercise}[$\star$]\label{exo:g8:cosine:5}
किसी \omterm{def:g6:shapes:triangles}{समकोण त्रिभुज} में $\theta$ की संलग्न भुजा $8$ नापती है और कर्ण
$10$। $\cos\theta$ निकालो, फिर $\theta$
($\cos^{-1}(0.8) \approx 36.9^\circ$)।
\end{exercise}

\begin{exercise}[$\star$]\label{exo:g8:cosine:6}
समझाओ कि किसी \omterm{def:g6:shapes:triangles}{समकोण त्रिभुज} के न्यून कोण के लिए $\cos\theta$ कभी
$1.2$ क्यों नहीं हो सकता।
\end{exercise}

\begin{exercise}[$\star\star$]\label{exo:g8:cosine:7}
$25$ m की एक ज़िप-लाइन एक मचान से ज़मीन तक उतरती है और आड़ी दिशा से
$12^\circ$ का कोण बनाती है ($\cos 12^\circ \approx 0.978$)। वह आड़ी
दूरी में कितनी लंबी है?
\end{exercise}

\begin{exercise}[$\star\star$]\label{exo:g8:cosine:8}
बिना कैलकुलेटर के क्रम में लगाओ: $\cos 20^\circ$, $\cos 70^\circ$,
$\cos 45^\circ$ (\cref{prop:g8:cosine:bounds} इस्तेमाल करो)। फिर
कैलकुलेटर से जाँचो।
\end{exercise}

\begin{exercise}[$\star\star$]\label{exo:g8:cosine:9}
$L$ लंबाई की एक सीढ़ी दीवार से टिकी है और सीढ़ी तथा \emph{ज़मीन} के
बीच का कोण $\theta$ है। उसका पैर दीवार से $1.2$ m दूर है और
$\theta = 68^\circ$ ($\cos 68^\circ \approx 0.375$)। $L$ निकालो,
फिर वह जितनी ऊँचाई तक पहुँचती है (पाइथागोरस से)।
\end{exercise}

\begin{exercise}[$\star\star$]\label{exo:g8:cosine:10}
$[BC]$ केंद्र $O$ और \omterm{def:g3:shapes:shapes}{त्रिज्या} $4.5$ cm वाले किसी वृत्त का \omterm{def:g4:geometry:circle}{व्यास} है,
और $A$ उस वृत्त का एक बिंदु है जिसके लिए $AB = 5.4$ cm है।
\begin{enumerate}
  \item $ABC$ का $A$ पर \omterm{def:g6:shapes:triangles}{समकोण} क्यों है?
  \item $AC$ निकालो, फिर $\cos \widehat B$।
\end{enumerate}
\end{exercise}

\begin{exercise}[$\star\star\star$]\label{exo:g8:cosine:11}
भुजा $1$ के \omterm{def:g6:shapes:triangles}{समबाहु त्रिभुज} को उसकी एक ऊँचाई के साथ काटकर \omterm{def:g3:division:half}{आधा} करो,
और उससे \cref{prop:g8:cosine:bounds} वाले पड़ाव
$\cos 60^\circ = \frac12$ का कारण दो। (ऊँचाई का पैर आधार पर कहाँ
गिरता है? पूरा त्रिकोणमितीय रूप \cref{ex:g9:trig:special} में है।)
\end{exercise}

\section{समस्या: यूक्लिड के संबंध, और पाइथागोरस फिर से}

\begin{problem}[{सप्ताहांत समस्या --- समकोण त्रिभुज का शीर्षलंब:
$AB^2 = BH \times BC$, $AH^2 = BH \times HC$, पाइथागोरस का दूसरा
प्रमाण, और वर्गमूल बनाने वाली एक मशीन}]
\label{pb:g8:cosine:1}
किसी \omterm{def:g6:shapes:triangles}{समकोण त्रिभुज} के \omterm{def:g6:shapes:triangles}{समकोण} से कर्ण पर \omterm{def:g6:lines:perp}{लंब} गिराओ: यह छोटा-सा
\omterm{def:g6:lines:objects}{रेखाखंड} --- \emph{शीर्षलंब} --- ऐसे संबंध मानता है जो इतने काम के
हैं कि यूक्लिड ने उन्हें अपनी \emph{एलिमेंट्स} के बीचोंबीच रखा।
इस समस्या में कोज्या का वही मूल गुण, ``अनुपात सिर्फ़ कोण पर टिका
है'' (\cref{def:g8:cosine:def}), उन सबको सिद्ध कर देता है; रास्ते
में तुम पाइथागोरस प्रमेय को बिलकुल अलग राह से फिर से सिद्ध करोगे,
और आख़िर में पैमाने-परकार की एक ऐसी मशीन बनाओगे जो हर पूर्ण संख्या
$n$ के लिए $\sqrt2$, $\sqrt5$, $\sqrt{n}$ बना देती है।

पूरे समय $ABC$ ऐसा त्रिभुज है जिसका $A$ पर \omterm{def:g6:shapes:triangles}{समकोण} है, और $H$, $A$ से
कर्ण $[BC]$ पर गिराए गए \omterm{def:g6:lines:perp}{लंब} का पैर है, यानी $H$, $B$ और $C$ के बीच
पड़ता है और $AH \perp BC$ है।

\medskip
\noindent\textbf{भाग I --- एक कोण, दो त्रिभुज।}
\begin{enumerate}
  \item $ABH$ और $ABC$ में त्रिभुज के कोणों के योग
        (\cref{thm:g7:triangles:sum}) से दिखाओ कि
        $\widehat{BAH} = \widehat{ACB}$: यानी शीर्षलंब \omterm{def:g6:shapes:triangles}{समकोण
        त्रिभुज} को ऐसे दो छोटे त्रिभुजों में काटता है जिनके कोण
        असली त्रिभुज \emph{जैसे ही} हैं।
  \item त्रिभुज $ABC$ ($A$ पर \omterm{def:g6:shapes:triangles}{समकोण}) और $ABH$ ($H$ पर \omterm{def:g6:shapes:triangles}{समकोण}) में
        कोण $\widehat B$ साझा है। दोनों में $\cos\widehat B$ लिखो,
        और \cref{def:g8:cosine:def} से निकालो कि
        \[
        AB^2 = BH \times BC .
        \]
        (बराबर अनुपातों से बराबर गुणनफलों तक जाने के लिए दोनों तरफ़
        दोनों हरों से गुणा करो,
        \cref{thm:g8:equations:balance}।)
  \item \omterm{def:g5:solids:def}{शीर्ष} $C$ पर वैसा ही जुड़वाँ संबंध लिखो और सिद्ध करो:
        \[
        AC^2 = CH \times CB .
        \]
  \item $3$--$4$--$5$ वाला त्रिभुज लो: $AB = 3$, $AC = 4$,
        $BC = 5$। सवाल 2 और 3 से $BH$ तथा $CH$ निकालो, और जाँचो कि
        $BH + HC = BC$।
  \item आम तौर पर, $BH$ और $CH$ को सिर्फ़ तीनों भुजाओं वाली भिन्नों
        के रूप में लिखो। किस अर्थ में पैर $H$ कर्ण को ``\omterm{def:g6:shapes:triangles}{समकोण}
        भुजाओं के वर्गों के \omterm{def:g7:prop:table}{समानुपात} में'' बाँटता है?
\end{enumerate}

\medskip
\noindent\textbf{भाग II --- पाइथागोरस फिर से, और शीर्षलंब।}
\begin{enumerate}[resume]
  \item सवाल 2 और 3 के संबंध जोड़ो और $BH + HC = BC$ इस्तेमाल करके
        निकालो
        \[
        AB^2 + AC^2 = BC^2 :
        \]
        यानी पाइथागोरस प्रमेय (\cref{thm:g8:pythagoras:direct}),
        फिर से सिद्ध --- और इस बार कहीं कोई \omterm{def:g6:measure:area}{क्षेत्रफल} वाली पहेली
        नहीं (\cref{exo:g8:pythagoras:11} से तुलना करो)।
  \item छोटे त्रिभुज $ABH$ में पाइथागोरस लगाकर
        $AH^2 = AB^2 - BH^2$ लिखो, $AB^2$ की जगह
        $BH \times BC$ रखो, और $BH$ बाहर निकालकर \emph{यूक्लिड का
        शीर्षलंब-संबंध} सिद्ध करो:
        \[
        AH^2 = BH \times HC .
        \]
  \item $ABC$ का \omterm{def:g6:measure:area}{क्षेत्रफल} दो तरह से निकालो --- \omterm{def:g6:shapes:triangles}{समकोण} भुजाओं को
        आधार और ऊँचाई मानकर, फिर कर्ण को आधार मानकर
        (\cref{thm:g7:areas:triangle}) --- और इससे तीसरा संबंध
        निकालो:
        \[
        AB \times AC = BC \times AH .
        \]
  \item फिर $3$--$4$--$5$ वाले त्रिभुज पर लौटो: सवाल 8 से $AH$
        निकालो, फिर सवाल 4 में मिले $BH$ तथा $CH$ के मानों से सवाल
        7 वाला शीर्षलंब-संबंध संख्याओं में जाँचो।
  \item किसी \omterm{def:g6:shapes:triangles}{समकोण त्रिभुज} में शीर्षलंब का पैर कर्ण को \omterm{def:g6:lines:objects}{रेखाखंडों}
        $BH = 2$ cm और $HC = 8$ cm में बाँटता है। शीर्षलंब $AH$
        निकालो।
\end{enumerate}

\medskip
\noindent\textbf{भाग III --- वर्गमूल बनाने वाली मशीन।}
एक \omterm{def:g6:lines:objects}{रेखाखंड} $[BC]$ बनाओ जो सिरे से सिरे जोड़े गए दो टुकड़ों से बना
हो: $BH = p$ और $HC = q$। \omterm{def:g4:geometry:circle}{व्यास} $[BC]$ वाला \omterm{def:g3:division:half}{आधा} वृत्त बनाओ, और मान
लो $A$ वह बिंदु है जहाँ $H$ पर $(BC)$ का \omterm{def:g6:lines:perp}{लंब} उस आधे वृत्त से मिलता
है।
\begin{enumerate}[resume]
  \item \cref{thm:g8:cosine:circle} से समझाओ कि त्रिभुज $ABC$ का
        $A$ पर \omterm{def:g6:shapes:triangles}{समकोण} क्यों है --- और इसलिए सवाल 7 लागू होता है तथा
        \[
        AH^2 = p \times q .
        \]
        (लंबाई $AH$ को $p$ और $q$ का \emph{गुणोत्तर
        माध्य}\index{गुणोत्तर माध्य} कहते हैं।)
  \item $p = 1$ और $q = 2$ लो। $AH^2$ क्या है?
        \cref{ex:g8:pythagoras:diagonal} वाली कौन-सी मशहूर लंबाई
        तुम्हारी पैमाने-परकार वाली आकृति ने अभी-अभी बना दी?
  \item वह नुस्ख़ा बताओ जो किसी भी पूर्ण संख्या $n \geq 1$ के लिए
        $\sqrt n$ लंबाई का \omterm{def:g6:lines:objects}{रेखाखंड} बना देता है, और बताओ कि $p$
        तथा $q$ के किस चुनाव से $\sqrt 5$ बनता है।
  \item मान लो $O$ आधे वृत्त का केंद्र है। समझाओ कि $AH$ \omterm{def:g3:shapes:shapes}{त्रिज्या}
        $\frac{p + q}{2}$ से कभी ज़्यादा क्यों नहीं हो सकता, और
        $H$ की किस जगह पर दोनों बराबर होते हैं।
  \item इससे यह असमिका निकालो: सभी धन संख्याओं $p$ और $q$ के लिए
        \[
        p \times q \leq \left(\frac{p + q}{2}\right)^2,
        \]
        और बराबरी ठीक तब जब $p = q$ हो। फिर उसे बिना किसी ज्यामिति
        के दोबारा सिद्ध करो: $\left(\frac{p+q}{2}\right)^2 - p q$
        फैलाओ और उसमें कोई उल्लेखनीय सर्वसमिका पहचानो
        (\cref{pb:g8:equations:1})।
\end{enumerate}
\end{problem}
